% Unit 082 English translation of chapter6.tex lines 1740--2123.
% Source: Methods of Algebra, Volume 2, commit
% 9a5803ff77dd3257484cb177f851a73770a59dd3 (CC BY 4.0).
% stable-unit-id: o014.aljabr2.chapter6.cup-product
% Coverage: Section 6.10 in full on cup products; stops immediately before
% Section 6.11 on Tate cohomology at line 2124.

% segment-id: o014.li2.chapter6.cup-product.h001
\section{The cup-product operation}\label{sec:cup-product}
\hypertarget{o014.aljabr2.chapter6.cup-product}{ }

% segment-id: o014.li2.chapter6.cup-product.p001
Let $G$ be an arbitrary group. As usual, write
$\otimes=\otimes_{\Bbbk}$; this notation denotes both the tensor product of
$\Bbbk$-modules and the tensor product of $G$-modules.

% segment-id: o014.li2.chapter6.cup-product.cv001
\begin{convention}
	Write $\Ext_G^\bullet:=\Ext_{\Bbbk[G]}^\bullet$ and
	$\cate{D}(G):=\cate{D}(G\dcate{Mod})$.
\end{convention}

% segment-id: o014.li2.chapter6.cup-product.p002
The aim of this section is to construct a cup-product operation on cohomology
for arbitrary $G$-modules $M_1$ and $M_2$,
% segment-id: o014.li2.chapter6.cup-product.q001
\[ \cup: \Hm^p(G, M_1) \otimes \Hm^q(G, M_2) \to \Hm^{p+q}(G, M_1 \otimes M_2), \quad p, q \in \Z_{\geq 0}, \]
which for $p=q=0$ gives the evident morphism
$M_1^G\otimes M_2^G\to(M_1\otimes M_2)^G$.

% segment-id: o014.li2.chapter6.cup-product.p003
Reformulated in terms of $\Ext_G^\bullet$, the aim is equivalently to define
% segment-id: o014.li2.chapter6.cup-product.q002
\[ \cup: \Ext^p_G(\Bbbk, M_1) \otimes \Ext^q_G(\Bbbk, M_2) \to \Ext_G^{p+q}(\Bbbk, M_1 \otimes M_2). \]

% segment-id: o014.li2.chapter6.cup-product.p004
Recall that
$\Ext^p_G(\Bbbk,M_1)\simeq\Hom_{\cate{D}(G)}(\Bbbk,M_1[p])$. Our first
viewpoint on the cup product is based on morphisms in the derived category.
Readers interested only in concrete computations with complexes may proceed
directly to Definition--Proposition~\ref{def:group-cup-3}.

% segment-id: o014.li2.chapter6.cup-product.d001
\begin{definition}
	Write the left-derived bifunctor of
	$\otimes:G\dcate{Mod}\times G\dcate{Mod}\to G\dcate{Mod}$ as
	$\otimesL$. This bifunctor gives $\cate{D}^-(G)$ the structure of a
	symmetric monoidal category, with the trivial $G$-module $\Bbbk$ as its
	unit object.
\end{definition}

% segment-id: o014.li2.chapter6.cup-product.p005
The tensor product $\otimes$ of complexes is defined, as usual, by taking the
total complex $\tot_{\oplus}$. The commutativity constraint in the symmetric
monoidal structure involves several signs; see
Proposition~\ref{prop:double-cplx-swap} for details.

% segment-id: o014.li2.chapter6.cup-product.p006
A projective $G$-module remains projective as a $\Bbbk$-module
(Proposition~\ref{prop:Ind-preservation}, with $H$ trivial), and is therefore
also flat. If $X$ and $Y$ are bounded-above complexes of $G$-modules, and
every term of $X$ is flat as a $\Bbbk$-module, then
% segment-id: o014.li2.chapter6.cup-product.q003
\[ X \otimes Y = X \otimesL Y, \quad Y \otimes X = Y \otimesL X; \]
see Proposition~\ref{prop:flat-K-flat}. In particular,
% segment-id: o014.li2.chapter6.cup-product.e001
\begin{equation}\label{eqn:G-mod-tensor-unit}
	\Bbbk \otimes Y = \Bbbk \otimesL Y \simeq Y \simeq Y \otimesL \Bbbk = Y \otimes \Bbbk.
\end{equation}

% segment-id: o014.li2.chapter6.cup-product.le001
\begin{lemma}\label{prop:group-pullback-monoidal}
	Let $\varphi:H\to G$ be a group homomorphism. Continue to write the functor
	on derived categories induced by the exact functor $\varphi^*$ as
	$\varphi^*:\cate{D}^-(G)\to\cate{D}^-(H)$. Then $\varphi^*$ is a
	monoidal functor with respect to $\otimesL$; see
	\cite[Definition 3.1.7]{Li1} and the remark following it. In other words,
	there are canonical isomorphisms satisfying all compatibility conditions,
% segment-id: o014.li2.chapter6.cup-product.q004
	\[ \varphi^*(\Bbbk) \simeq \Bbbk, \quad \varphi^* X_1 \otimesL \varphi^* X_2 \simeq \varphi^* (X_1 \otimesL X_2). \]
\end{lemma}

% segment-id: o014.li2.chapter6.cup-product.pf001
\begin{proof}
	The isomorphism $\varphi^*\Bbbk\simeq\Bbbk$ (indeed, an equality) and all
	its compatibilities are immediate. Furthermore, if a $G$-module $M$ is
	flat as a $\Bbbk$-module, then so is $\varphi^*M$. Since $\otimesL$ can be
	computed using resolutions that are flat over $\Bbbk$, the canonical
	isomorphism
	$\varphi^*X_1\otimesL\varphi^*X_2\simeq\varphi^*(X_1\otimesL X_2)$
	is obtained by reducing to the case in which $X_1$ and $X_2$ are termwise
	flat over $\Bbbk$.
\end{proof}

% segment-id: o014.li2.chapter6.cup-product.p007
We can now formulate a preliminary version of the cup product in the derived
category; the operation simply reflects the tensor-product operation. For the
moment, write the isomorphism
$\Bbbk\rightiso\Bbbk\otimes\Bbbk=\Bbbk\otimesL\Bbbk$ given by
\eqref{eqn:G-mod-tensor-unit} as $\delta$.

% segment-id: o014.li2.chapter6.cup-product.d002
\begin{definition}\label{def:group-cup-1}
	Let $X_1$ and $X_2$ be objects of $\cate{D}^-(G)$. Define the
	$\Bbbk$-module homomorphism
% segment-id: o014.li2.chapter6.cup-product.q005
	\[ \overset{\mathrm{L}}{\cup}: \Hom_{\cate{D}(G)}(\Bbbk, X_1) \otimes \Hom_{\cate{D}(G)}(\Bbbk, X_2) \to \Hom_{\cate{D}(G)}\left(\Bbbk, X_1 \otimesL X_2\right) \]
	as follows: $\alpha\overset{\mathrm{L}}{\cup}\beta$ is the composite
	morphism
% segment-id: o014.li2.chapter6.cup-product.q006
	\[ \Bbbk \xrightarrow{\delta} \Bbbk \otimes \Bbbk = \Bbbk \otimesL \Bbbk \xrightarrow{\alpha \otimesL \beta} X_1 \otimesL X_2. \]
\end{definition}

% segment-id: o014.li2.chapter6.cup-product.p008
Clearly, the map
$(\alpha,\beta)\mapsto\alpha\overset{\mathrm{L}}{\cup}\beta$ is bilinear,
so the definition is valid. The operation can also be understood as the
composition of morphisms. Indeed, from the general properties of monoidal
categories, it is easy to see that
$\alpha\overset{\mathrm{L}}{\cup}\beta$ is given by the following composite:
% segment-id: o014.li2.chapter6.cup-product.g001
\[\begin{tikzcd}[row sep=small]
	\Hom_{\cate{D}(G)}(\Bbbk, X_1) \otimes \Hom_{\cate{D}(G)}(\Bbbk, X_2) \arrow[d] \\
	\Hom_{\cate{D}(G)}\left( \Bbbk \otimes X_2, X_1 \otimesL X_2\right) \otimes \Hom_{\cate{D}(G)}(\Bbbk \otimes \Bbbk, \Bbbk \otimes X_2) \arrow[d] \\
	\Hom_{\cate{D}(G)}\left(\Bbbk, X_1 \otimesL X_2 \right)
\end{tikzcd}\]
The first part is based on \eqref{eqn:G-mod-tensor-unit}, while the second is
composition of morphisms in $\cate{D}(G)$.

% segment-id: o014.li2.chapter6.cup-product.p009
The following properties are formal consequences of
Definition~\ref{def:group-cup-1}, or of the reformulation above; throughout,
$X_i$ and $X$ are objects of $\cate{D}^-(G)$.

% segment-id: o014.li2.chapter6.cup-product.l001
\begin{description}
% segment-id: o014.li2.chapter6.cup-product.i001
	\item[Functoriality] The operation $\overset{\mathrm{L}}{\cup}$ is
	functorial in $X_1$ and $X_2$.
% segment-id: o014.li2.chapter6.cup-product.i002
	\item[Unit] Under the isomorphisms in
	\eqref{eqn:G-mod-tensor-unit}, both the left and the right action of
	$\identity\in\End_{\cate{D}(G)}(\Bbbk)$ through
	$\overset{\mathrm{L}}{\cup}$ on $\Hom_{\cate{D}(G)}(\Bbbk,X)$ are the
	identity.
% segment-id: o014.li2.chapter6.cup-product.i003
	\item[Associativity] Let
	$\alpha_i\in\Hom_{\cate{D}(G)}(\Bbbk,X_i)$ for $i=1,2,3$. Up to the
	associativity constraint for $\otimesL$,
% segment-id: o014.li2.chapter6.cup-product.q007
	\[ \alpha_1 \overset{\mathrm{L}}{\cup} (\alpha_2 \overset{\mathrm{L}}{\cup} \alpha_3) = (\alpha_1 \overset{\mathrm{L}}{\cup} \alpha_2) \overset{\mathrm{L}}{\cup} \alpha_3 . \]
% segment-id: o014.li2.chapter6.cup-product.i004
	\item[Commutativity] Under the commutativity constraint
	$X_1\otimesL X_2\rightiso X_2\otimesL X_1$, the element
	$\alpha\overset{\mathrm{L}}{\cup}\beta$ is mapped to
	$\beta\overset{\mathrm{L}}{\cup}\alpha$. To see this, it suffices to
	apply Definition~\ref{def:group-cup-1} and the following commutative
	diagram:
% segment-id: o014.li2.chapter6.cup-product.g002
	\[\begin{tikzcd}
		\Bbbk \arrow[r, "\delta"] \arrow[rd, "\delta"'] & \Bbbk \otimesL \Bbbk \arrow[d] \arrow[r, "{\alpha \otimesL \beta}"] & X_1 \otimesL X_2 \arrow[d] \\
		& \Bbbk \otimesL \Bbbk \arrow[r, "{\beta \otimesL \alpha}"] & X_2 \otimesL X_1
	\end{tikzcd}\]
	All vertical arrows are commutativity constraints for $\otimesL$. This is
	a general fact about symmetric monoidal categories.
\end{description}

% segment-id: o014.li2.chapter6.cup-product.p010
We next define the cup product at the level of cohomology of $G$-modules;
several signs will be involved.

% segment-id: o014.li2.chapter6.cup-product.d003
\begin{definition}[Cup product]\label{def:group-cup-2}
	\index{cup product}
	\index[sym1]{cup@$\cup$}
	Let $M_1$ and $M_2$ be $G$-modules and let $p_1,p_2\in\Z$. In
	Definition~\ref{def:group-cup-1}, take $X_i=M_i[p_i]$, and then compose
	$\overset{\mathrm{L}}{\cup}$ with the canonical morphisms
% segment-id: o014.li2.chapter6.cup-product.q008
	\begin{align*}
		M_1[p_1] \otimesL M_2[p_2] & \xrightarrow{\text{can}} M_1[p_1] \otimes M_2[p_2] \\
		& \xrightarrow[\sim]{\theta'} (M_1[p_1] \otimes M_2) [p_2] \xrightarrow[\sim]{\theta} (M_1 \otimes M_2)[p_1 + p_2]
	\end{align*}
	in $\cate{D}(G)$, where
% segment-id: o014.li2.chapter6.cup-product.l002
	\begin{compactitem}
% segment-id: o014.li2.chapter6.cup-product.i005
		\item $\mathrm{can}$ is the canonical morphism associated with the
		left-derived bifunctor;
% segment-id: o014.li2.chapter6.cup-product.i006
		\item the canonical isomorphisms $\theta$ and $\theta'$ are defined as
		in Proposition~\ref{prop:tot-shift}.
	\end{compactitem}
	This yields the cup product on cohomology,
% segment-id: o014.li2.chapter6.cup-product.q009
	\[ \cup: \Ext^{p_1}_G(\Bbbk, M_1) \otimes \Ext^{p_2}_G(\Bbbk, M_2) \to \Ext^{p_1 + p_2}_G(\Bbbk, M_1 \otimes M_2), \]
	or, in other notation,
	$\cup:\Hm^{p_1}(G,M_1)\otimes\Hm^{p_2}(G,M_2)
	\to\Hm^{p_1+p_2}(G,M_1\otimes M_2)$.
\end{definition}

% segment-id: o014.li2.chapter6.cup-product.p011
If $p_1=p_2=0$, the cup product is obtained by taking the tensor product of
two morphisms $\Bbbk\to M_1$ and $\Bbbk\to M_2$ in $G\dcate{Mod}$, namely
$\Bbbk\simeq\Bbbk\otimes\Bbbk\to M_1\otimes M_2$. The result is precisely
the evident map $(M_1)^G\otimes(M_2)^G\to(M_1\otimes M_2)^G$.

% segment-id: o014.li2.chapter6.cup-product.p012
Inspecting the definition, the isomorphism
$M_1[p_1]\otimes M_2[p_2]\rightiso(M_1\otimes M_2)[p_1+p_2]$ involved has
the following concrete description: on its only nonzero term
$M_1[p_1]^{-p_1}\otimes M_2[p_2]^{-p_2}=M_1\otimes M_2$, it acts as
$(-1)^{p_1p_2}$; the sign comes from $\theta'$.

% segment-id: o014.li2.chapter6.cup-product.p013
In addition to the functoriality and unit law already formulated in the
derived category, the cup product has the following properties.

% segment-id: o014.li2.chapter6.cup-product.l003
\begin{description}
% segment-id: o014.li2.chapter6.cup-product.i007
	\item[Associativity] For $\alpha_i\in\Hm^{p_i}(G,M_i)$ with
	$i=1,2,3$, one has
% segment-id: o014.li2.chapter6.cup-product.q010
	\[ (\alpha_1 \cup \alpha_2) \cup \alpha_3 = \alpha_1 \cup (\alpha_2 \cup \alpha_3). \]
	This property follows easily from the associativity of
	$\overset{\mathrm{L}}{\cup}$, but the sign mentioned above still needs to
	be taken into account. In fact, the sign on both sides is
	$(-1)^{p_1p_2+p_2p_3+p_1p_3}$.

% segment-id: o014.li2.chapter6.cup-product.i008
	\item[Graded commutativity] Let
	$c:M_1\otimes M_2\rightiso M_2\otimes M_1$ be the isomorphism determined
	by $x\otimes y\mapsto y\otimes x$. For all
	$\alpha\in\Hm^p(G,M_1)$ and $\beta\in\Hm^q(G,M_2)$, one has
% segment-id: o014.li2.chapter6.cup-product.q011
	\[ \Hm^{p+q}(c)(\alpha \cup \beta) = (-1)^{pq} \beta \cup \alpha. \]

	Indeed, by the commutativity of $\overset{\mathrm{L}}{\cup}$, it suffices
	to check that the following diagram commutes:
% segment-id: o014.li2.chapter6.cup-product.g003
	\[\begin{tikzcd}
		{M_1[p]} \otimesL {M_2[q]} \arrow[d] \arrow[r, "\text{can}"] & {M_1[p]} \otimes {M_2[q]} \arrow[d] \arrow[r, "\text{extract $[q]$ first}" inner sep=0.6em] & (M_1 \otimes M_2){[p+q]} \arrow[d, "{(-1)^{pq} c[p+q]}"] \\
		{M_2[q]} \otimesL {M_1[p]} \arrow[r, "\text{can}"'] & {M_2[q]} \otimes {M_1[p]} \arrow[r, "\text{extract $[p]$ first}"' inner sep=0.6em] & (M_2 \otimes M_1) {[p+q]}
	\end{tikzcd}\]
	The vertical arrows on the left and in the middle are the commutativity
	constraints for tensor products in the derived category or the category of
	complexes. The left square plainly commutes. For the right square, the last
	part of Proposition~\ref{prop:tot-shift} says that the following diagram
	commutes:
% segment-id: o014.li2.chapter6.cup-product.g004
	\[\begin{tikzcd}[column sep=large]
		{M_1[p]} \otimes {M_2[q]} \arrow[d] \arrow[r, "\text{extract $[q]$ first}"] & (M_1 \otimes M_2) {[p+q]} \arrow[d, "{c[p+q]}"] \\
		{M_2[q]} \otimes {M_1[p]} \arrow[r, "\text{extract $[q]$ first}"'] & (M_2 \otimes M_1) {[p+q]}.
	\end{tikzcd}\]
	However, by the anticommutative diagram in
	Proposition~\ref{prop:tot-shift} and its straightforward generalization,
	the isomorphism
	$M_2[q]\otimes M_1[p]\rightiso(M_2\otimes M_1)[p+q]$ obtained by
	extracting $[q]$ first differs by $(-1)^{pq}$ from the one obtained by
	extracting $[p]$ first. This proves commutativity of the right square in
	the original diagram.

% segment-id: o014.li2.chapter6.cup-product.i009
	\item[Cohomology algebra] In particular,
	$\Hm^\bullet(G,\Bbbk):=\bigoplus_p\Hm^p(G,\Bbbk)$ becomes a graded
	$\Bbbk$-algebra under $\cup$, and its multiplication satisfies the graded
	commutativity law above,
	$\beta\cup\alpha=(-1)^{pq}\alpha\cup\beta$. In fact, one can prove that
	$\Hm^\bullet(G,\Bbbk)$ is the same as the graded $\Bbbk$-algebra in
	Definition~\ref{def:Ext-algebra},
% segment-id: o014.li2.chapter6.cup-product.q012
	\[ \Ext_G(\Bbbk) := \bigoplus_p \Ext^p_G(\Bbbk, \Bbbk) = \bigoplus_p \Hm^p(G, \Bbbk). \]
	% Reference: Brown, Cohomology of Groups, V.4.6 Theorem.

% segment-id: o014.li2.chapter6.cup-product.i010
	\item[Bimodule structure]
	Under the operation $\cup$, $\bigoplus_p\Hm^p(G,M)$ becomes a graded
	bimodule with left and right actions by $\Hm^\bullet(G,\Bbbk)$.
\end{description}

% segment-id: o014.li2.chapter6.cup-product.p014
A group homomorphism $\varphi:H\to G$ induces a homomorphism of cohomology
algebras $\varphi^*:\Hm^\bullet(G,\Bbbk)\to\Hm^\bullet(H,\Bbbk)$. This is an
immediate consequence of the following result.

% segment-id: o014.li2.chapter6.cup-product.pr001
\begin{proposition}[Change of groups]
	Use the notation of Definition~\ref{def:group-cup-2}. Let
	$\varphi:H\to G$ be a group homomorphism. Write
	$\varphi^*:\Hm^{p_i}(G,M_i)\to\Hm^{p_i}(H,\varphi^*M_i)$ for the
	canonical homomorphism in \eqref{eqn:change-of-group-coh}, for $i=1,2$.
	Then
% segment-id: o014.li2.chapter6.cup-product.q013
	\[ \varphi^*(\alpha \cup \beta) = \varphi^* \alpha \cup \varphi^* \beta. \]

	As a special case, if $H\subset G$ is a subgroup, the restriction map
	$\mathrm{res}^n:\Hm^n(G,M)\to\Hm^n(H,M)$ discussed in
	\S\ref{sec:group-finite-index} preserves cup products.
\end{proposition}

% segment-id: o014.li2.chapter6.cup-product.pf002
\begin{proof}
	Apply $\varphi^*$ to every morphism in
	Definition~\ref{def:group-cup-1}, and then use the fact that $\varphi^*$ is
	a monoidal functor (Lemma~\ref{prop:group-pullback-monoidal}). In
	$\Ext_H^{p_1+p_2}(\Bbbk,\varphi^*M_1\otimes\varphi^*M_2)$, this gives
% segment-id: o014.li2.chapter6.cup-product.g005
	\begin{multline*}
		\varphi^*(\alpha \cup \beta) \xlongequal{\text{identified with}} \;\text{the composite morphism} \\
		\left[ \Bbbk \xrightarrow{\delta} \Bbbk \otimesL \Bbbk \xrightarrow{\varphi^* \alpha \otimesL \varphi^* \beta} \varphi^* M_1[p_1] \otimesL \varphi^* M_2[p_2] \to (\varphi^* M_1 \otimes \varphi^* M_2)[p_1 + p_2] \right].
	\end{multline*}
	Here,
	$\varphi^*\alpha\in\Ext_H^{p_1}(\Bbbk,\varphi^*M_1)
	\simeq\Hom_H(\varphi^*\Bbbk,\varphi^*M_1[p_1])$ is obtained by applying
	the functor $\varphi^*$ to the morphism $\alpha$. By
	Remark~\ref{rem:change-of-group-Yoneda}, however, this element is also the
	image of $\alpha$ under the canonical homomorphism
	\eqref{eqn:change-of-group-coh}; the same applies to $\varphi^*\beta$.
	The result is $\varphi^*\alpha\cup\varphi^*\beta$.
\end{proof}

% segment-id: o014.li2.chapter6.cup-product.p015
Our next aim is to describe the cup product through the standard complex
$C(G,M)$ of Proposition~\ref{prop:groupcoh-cochain}. This description is not
only more concrete; it also lifts the cup product canonically to the level of
complexes. We begin with the definition. Later,
Proposition~\ref{prop:std-cplx-cup} will show that it agrees with the previous
version.

% segment-id: o014.li2.chapter6.cup-product.dp001
\begin{definition-proposition}[Cup product on the standard complex]\label{def:group-cup-3}
	\index{cup product}
	Let $M_1$ and $M_2$ be $G$-modules. For all $p_1,p_2\in\Z$, define the
	homomorphisms
% segment-id: o014.li2.chapter6.cup-product.g006
	\[\begin{tikzcd}[row sep=small, column sep=small]
		\cup: C^{p_1}(G, M_1) \otimes C^{p_2}(G, M_2) \arrow[r] & C^{p_1 + p_2}(G, M_1 \otimes M_2) \\
		f_1 \otimes f_2 \arrow[mapsto, r] & f_1(g_1, \ldots, g_{p_1}) \otimes (g_1 \cdots g_{p_1} f_2(g_{p_1 + 1}, \ldots, g_{p_1 + p_2}))
	\end{tikzcd}\]
	where $g_1,\ldots,g_{p_1+p_2}\in G$.
% segment-id: o014.li2.chapter6.cup-product.l004
	\begin{enumerate}[(i)]
% segment-id: o014.li2.chapter6.cup-product.i011
		\item These homomorphisms satisfy associativity (taking $M_1$, $M_2$,
		and $M_3$ into account) and the Leibniz rule
% segment-id: o014.li2.chapter6.cup-product.q014
		\[ d(f_1 \cup f_2) = df_1 \cup f_2 + (-1)^{p_1} f_1 \cup d f_2; \]
		in particular, $\cup$ gives a morphism of complexes
		$\cup:C(G,M_1)\otimes C(G,M_2)\to C(G,M_1\otimes M_2)$.
% segment-id: o014.li2.chapter6.cup-product.i012
		\item The normalized standard complex introduced in
		Remark~\ref{rem:nor-cochain} is closed under $\cup$, and hence there is
		a morphism
% segment-id: o014.li2.chapter6.cup-product.q015
		\[ \cup: \overline{C}(G, M_1) \otimes \overline{C}(G, M_2) \to \overline{C}(G, M_1 \otimes M_2). \]
	\end{enumerate}
\end{definition-proposition}

% segment-id: o014.li2.chapter6.cup-product.pf003
\begin{proof}
	Associativity follows immediately from the formula for $\cup$, while the
	Leibniz rule is also easy to verify. Let
% segment-id: o014.li2.chapter6.cup-product.q016
	\[ \varphi \in C^p(G, M_1), \quad \psi \in C^q(G, M_2). \]
	In
% segment-id: o014.li2.chapter6.cup-product.g007
	\begin{multline*}
		(d\varphi \cup \psi)(g_1, \ldots, g_{p+q+1}) = \\
		\left( g_1 \varphi(g_2, \ldots, g_{p+1}) + \sum_{k=1}^p (-1)^k \varphi(\ldots, g_k g_{k+1}, \ldots) + (-1)^{p+1} \varphi(g_1, \ldots, g_p) \right) \\
		\otimes (g_1 \cdots g_{p+1}) \psi(g_{p+2}, \ldots, g_{p+q+1})
	\end{multline*}
and
% segment-id: o014.li2.chapter6.cup-product.g008
	\begin{multline*}
		(-1)^p (\varphi \cup d\psi)(g_1, \ldots, g_{p+q+1}) = \varphi(g_1, \ldots, g_p) \otimes \\
		(g_1 \cdots g_p) \bigg( (-1)^p g_{p+1} \psi(g_{p+2}, \ldots, g_{p+q+1}) + \sum_{k=p+1}^{p+q} (-1)^k \psi(\ldots, g_k g_{k+1}, \ldots) \\
		+ (-1)^{p+q+1} \psi(g_{p+1}, \ldots, g_{p+q}) \bigg),
	\end{multline*}
	the term
	$\varphi(g_1,\ldots,g_p)\otimes(g_1\cdots g_{p+1})
	\psi(g_{p+2},\ldots)$ occurs respectively as the last and first term,
	and the two cancel. The sum of the two expressions is therefore
	$d(\varphi\cup\psi)(g_1,\ldots,g_{p+q+1})$.

	Notice that the Leibniz rule is equivalent to saying that, as $p_1,p_2$
	vary, $\cup$ gives a morphism of complexes. This proves~(i).

	% Editorial correction: the stated vanishing on every input position
	% requires both cochains to be normalized; “or” in the source is insufficient.
	If both $f_1$ and $f_2$ belong to the normalized standard complex, the
	definition implies that
	$(f_1\cup f_2)(g_1,\ldots,g_{p_1+p_2})=0$ whenever one of
	$g_1,\ldots,g_{p_1+p_2}$ equals $1_G$. This also proves~(ii).
\end{proof}

% segment-id: o014.li2.chapter6.cup-product.p016
Recall that the standard complex $C(G,M)$ essentially computes
$\RHom_G(\Bbbk,M)=\Hom_G^\bullet(\mathsf{L},M)$ through the free resolution
$\mathsf{L}\to\Bbbk$, where $M$ is a $G$-module. More precisely, for
$\mathsf{L}$ we use the chain complex $(\mathsf{L}_n,\partial_n)_n$ in
Definition~\ref{def:resolution-trivial-module}. By setting
$\mathsf{L}^n:=\mathsf{L}_{-n}$ while retaining the morphisms $\partial_n$,
$\mathsf{L}$ may be regarded as a complex.

% segment-id: o014.li2.chapter6.cup-product.le002
\begin{lemma}\label{prop:cup-resolution}
	Write the free resolution of Definition~\ref{def:resolution-trivial-module}
	as $\epsilon:\mathsf{L}\to\Bbbk$. Then:
% segment-id: o014.li2.chapter6.cup-product.l005
	\begin{enumerate}[(i)]
% segment-id: o014.li2.chapter6.cup-product.i013
		\item $\epsilon\otimes\epsilon:\mathsf{L}\otimes\mathsf{L}
		\to\Bbbk\otimes\Bbbk\simeq\Bbbk$ is a projective resolution of the
		$G$-module $\Bbbk$;
% segment-id: o014.li2.chapter6.cup-product.i014
		\item one can define a morphism of complexes
		$\Delta:\mathsf{L}\to\mathsf{L}\otimes\mathsf{L}$ whose degree-$-n$
		term is given by
% segment-id: o014.li2.chapter6.cup-product.g009
		\[\begin{tikzcd}[row sep=tiny]
			\Delta_n: \mathsf{L}_n \arrow[r] & \bigoplus_{p=0}^n \mathsf{L}_p \otimes \mathsf{L}_{n-p} \\
			(g_0, \ldots, g_n) \arrow[mapsto, r] & \left( (-1)^{p(n-p)} (g_0, \ldots, g_p) \otimes (g_p, \ldots, g_n) \right)_{p=0}^n
		\end{tikzcd}\]
		and which makes the following diagram commute:
% segment-id: o014.li2.chapter6.cup-product.g010
		\[\begin{tikzcd}
			\mathsf{L} \arrow[r, "\Delta"] \arrow[d, "\epsilon"'] & \mathsf{L} \otimes \mathsf{L} \arrow[d, "{\epsilon \otimes \epsilon}"] \\
			\Bbbk \arrow[r, "\sim", "\delta"'] & \Bbbk \otimes \Bbbk .
		\end{tikzcd}\]
	\end{enumerate}
\end{lemma}

% segment-id: o014.li2.chapter6.cup-product.pf004
\begin{proof}
	For~(i), notice that every term of $\mathsf{L}\otimes\mathsf{L}$ remains a
	projective $G$-module; one may use
	Proposition~\ref{prop:group-tensor-proj}. To prove that
	$\epsilon\otimes\epsilon$ is a quasi-isomorphism, factor the morphism as
% segment-id: o014.li2.chapter6.cup-product.q017
	\[ \mathsf{L} \otimes \mathsf{L} \to \mathsf{L} \otimes \Bbbk \to \Bbbk \otimes \Bbbk. \]
	Since $\mathsf{L}$ consists of flat $\Bbbk$-modules, both parts are
	quasi-isomorphisms.

	For~(ii), it is clear that every $\Delta_n$ is a homomorphism of
	$G$-modules. A direct calculation, split into the various cases, verifies
	that $\Delta$ is indeed a morphism of complexes. Commutativity of the
	diagram is immediate.
\end{proof}

% segment-id: o014.li2.chapter6.cup-product.p017
In comparison with the cup product in topology, $\Delta$ here is analogous
to the ``diagonal embedding'' of a space. This morphism can also be
interpreted through the Alexander--Whitney map to be introduced in
\S\ref{sec:bisimplicial}; see the exercises in \CHref{sec:simplicial}.

% segment-id: o014.li2.chapter6.cup-product.p018
Notice that, because $\epsilon\otimes\epsilon$ is a quasi-isomorphism, the
general theory of complexes guarantees the existence of a morphism
$\mathsf{L}\to\mathsf{L}\otimes\mathsf{L}$ that makes the diagram in~(ii)
commute, and this morphism is unique up to homotopy
(Theorem~\ref{prop:resolution-homotopy}). The value of
Lemma~\ref{prop:cup-resolution}~(ii) lies in its explicit formula.

% segment-id: o014.li2.chapter6.cup-product.p019
Return to the description of $\overset{\mathrm{L}}{\cup}$ in
Definition~\ref{def:group-cup-1}. Represent the morphisms $\alpha$ and
$\beta$ in $\cate{D}(G)$ by morphisms of complexes
$\mathsf{L}\to X_1$ and $\mathsf{L}\to X_2$, respectively, where the $X_i$
are bounded-above complexes, and lift $\delta$ to $\Delta$ through
Lemma~\ref{prop:cup-resolution}. The composite of
$\alpha\overset{\mathrm{L}}{\cup}\beta
\in\Hom_{\cate{D}(G)}(\Bbbk,X_1\otimesL X_2)$ with the canonical morphism
$X_1\otimesL X_2\xrightarrow{\text{can}}X_1\otimes X_2$ is given concretely
by the composite of morphisms of complexes
% segment-id: o014.li2.chapter6.cup-product.g011
\[\begin{tikzcd}[row sep=small]
	\mathsf{L} \arrow[r, "\Delta"] & \mathsf{L} \otimes \mathsf{L} \arrow[r, "{\alpha \otimes \beta}"] & X_1 \otimes X_2. \\
	& \mathsf{L} \otimesL \mathsf{L} \arrow[equal, u] &
\end{tikzcd}\]

% segment-id: o014.li2.chapter6.cup-product.p020
For $G$-modules $M_i$ and $p_i\in\Z$, one has
$\Hm^{p_i}\Hom^\bullet(\mathsf{L},M_i)\simeq\Hm^{p_i}(G,M_i)$, for
$i=1,2$.

% segment-id: o014.li2.chapter6.cup-product.p021
We next interpret the cup product through $\Hom$ complexes. For arbitrary
complexes of $G$-modules $X,Y,X',Y'$ and $m,n\in\Z$, define the
$\Bbbk$-module homomorphism
% segment-id: o014.li2.chapter6.cup-product.e002
\begin{equation}\label{eqn:group-cup-3-aux}
	\mu^{m, n}: \Hom_G^m(X, Y) \otimes \Hom_G^n(X', Y') \to \Hom_G^{m+n}(X \otimes X', Y \otimes Y'),
\end{equation}
as follows. For an element $f\otimes g$ on the left, define
% segment-id: o014.li2.chapter6.cup-product.g012
\begin{gather*}
	\mu^{m, n}(f \otimes g)(x \otimes x') := (-1)^{pn} f(x) \otimes g(x') \; \in Y^{p+m} \otimes (Y')^{q+n}, \\
	x \otimes x' \in X^p \otimes (X')^q .
\end{gather*}
A direct verification shows that all the $\mu^{m,n}$ give a morphism of
complexes
% segment-id: o014.li2.chapter6.cup-product.q018
\[ \mu: \Hom_G^\bullet(X, Y) \otimes \Hom_G^\bullet(X', Y') \to \Hom_G^\bullet(X \otimes X', Y \otimes Y'), \]
or, in other words, that \eqref{eqn:group-cup-3-aux} satisfies the Leibniz
rule. The sign $(-1)^{pn}$ in the definition can be explained by the Koszul
sign rule to be discussed in \S\ref{sec:alg-in-monoidal-cat} and the
exercises in \CHref{sec:monads}, since the formula interchanges the positions
of $g$ and $x$.

% segment-id: o014.li2.chapter6.cup-product.le003
\begin{lemma}\label{prop:std-cplx-cup-aux}
	Apply the operation above with $X=X'=\mathsf{L}$, $Y=M_1$, and $Y'=M_2$.
	Then the morphism of complexes
% segment-id: o014.li2.chapter6.cup-product.e003
	\begin{multline}\label{eqn:group-cup-3-aux1}
		\Hom^\bullet_G(\mathsf{L}, M_1) \otimes \Hom^\bullet_G(\mathsf{L}, M_2) \xrightarrow{\mu} \Hom^\bullet_G(\mathsf{L} \otimes \mathsf{L}, M_1 \otimes M_2) \\
		\xrightarrow{\Delta^*} \Hom^\bullet(\mathsf{L}, M_1 \otimes M_2)
	\end{multline}
	induces on cohomology the cup product of
	Definition~\ref{def:group-cup-2}.
\end{lemma}

% segment-id: o014.li2.chapter6.cup-product.pf005
\begin{proof}
	For $i=1,2$, use Lemma~\ref{prop:Hom-cplx-d-translate} to identify
	$\Hom_G^{p_i}(\mathsf{L},M_i)$ with
	$\Hom_G^0(\mathsf{L},M_i[p_i])$. Choose $f_i$ in it satisfying
	$d_{\Hom^\bullet}f_i=0$, and consider the image of $f_1\otimes f_2$ under
\eqref{eqn:group-cup-3-aux}. Since $f_i$ can be nonzero only on
$\mathsf{L}_{p_i}$, the sign in this image is $(-1)^{p_1p_2}$, exactly the
sign in Definition~\ref{def:group-cup-2}. The remaining verification is
immediate.
\end{proof}

% segment-id: o014.li2.chapter6.cup-product.p022
Recall that $\Hom^\bullet(\mathsf{L},M_i)\simeq C(G,M_i)$, so that
$\Hm^{p_i}(G,M_i)\simeq\Hm^{p_i}(C(G,M_i))$. The concrete map is given in
the proof of Proposition~\ref{prop:groupcoh-cochain}, for $i=1,2$.

% segment-id: o014.li2.chapter6.cup-product.pr002
\begin{proposition}\label{prop:std-cplx-cup}
	The cup product in Definition~\ref{def:group-cup-2} is the same as the
	composite homomorphism
% segment-id: o014.li2.chapter6.cup-product.g013
	\begin{multline*}
		\Hm^{p_1}(G, M_1) \otimes \Hm^{p_2}(G, M_2) \simeq \Hm^{p_1}(C(G, M_1)) \otimes \Hm^{p_2}(C(G, M_2)) \\
		\xrightarrow{\kappa} \Hm^{p_1 + p_2}\left( C(G, M_1) \otimes C(G, M_2) \right) \\
		\xrightarrow{\Hm^{p_1 + p_2}(\cup)} \Hm^{p_1 + p_2}(C(G, M_1 \otimes M_2)) \simeq \Hm^{p_1 + p_2}(G, M_1 \otimes M_2),
	\end{multline*}
	where the canonical morphism $\kappa$ is as described immediately before
	Theorem~\ref{prop:Kunneth-homology}. The conclusion is the same when the
	normalized standard complex is used.
\end{proposition}

% segment-id: o014.li2.chapter6.cup-product.pf006
\begin{proof}
	Interpret the cup product as in Lemma~\ref{prop:std-cplx-cup-aux}. Take
	$f_i\in C^{p_i}(G,M_i)$. For the bar notation in the display below, put
	$p:=p_1$. For $g_1,\ldots,g_{p_1+p_2}\in G$, the element
% segment-id: o014.li2.chapter6.cup-product.q019
	\[ (g_1 | \cdots | g_{p_1 + p_2}) := \left(1_G, g_{[1, 1]}, \ldots, g_{[1, p_1 + p_2]} \right) \in \mathsf{L}_{p_1 + p_2} \]
	has, after applying $\Delta$, the following component in
	$\mathsf{L}_{p_1}\otimes\mathsf{L}_{p_2}$:
% segment-id: o014.li2.chapter6.cup-product.g014
	\begin{multline*}
		(-1)^{p_1 p_2} \left( 1_G, \ldots, g_{[1, p_1]}\right) \otimes \left(g_{[1, p_1]}, \ldots, g_{[1, p_1 + p_2]} \right) \\
		= (-1)^{p_1 p_2} (g_1 | \cdots | g_p) \otimes g_1 \cdots g_{p_1} \left( g_{p_1 + 1} |\cdots| g_{p_1 + p_2} \right).
	\end{multline*}

	Under $\mu^{p_1,p_2}(f_1\otimes f_2)$ in
	\eqref{eqn:group-cup-3-aux}, this element is mapped to
% segment-id: o014.li2.chapter6.cup-product.q020
	\[ f_1(g_1, \ldots, g_{p_1}) \otimes g_1 \cdots g_{p_1} f_2(g_{p_1 + 1}, \cdots, g_{p_1 + p_2}). \]
	Thus the operation in \eqref{eqn:group-cup-3-aux1} is indeed compatible
	with $\cup$ on $C(G,M)$. The version for the normalized standard complex
	is the same.
\end{proof}

% segment-id: o014.li2.chapter6.cup-product.p023
Notice that although the description of the cup product on the standard
complex is direct and transparent, graded commutativity is difficult to
derive directly from it; in the framework of the standard complex, a rather
roundabout proof is required. This book does not take that route.

% segment-id: o014.li2.chapter6.cup-product.r001
\begin{remark}
	The structure exhibited by
	Definition--Proposition~\ref{def:group-cup-3} is genuinely richer than its
	cohomological version in Definition~\ref{def:group-cup-2}, because it
	equips the entire standard complex with a differential graded structure.
	Thus $C(G,\Bbbk)$ under $\cup$ becomes a differential graded algebra in
	the sense of Definition~\ref{def:dg-algebra-preview} or
	\ref{def:dg-algebra}; the same holds for $\overline{C}(G,\Bbbk)$. The
	cohomological version of the cup product is obtained only after taking
	$\Hm^\bullet$: it is graded, but not differential.
\end{remark}

% segment-id: o014.li2.chapter6.cup-product.p024
The cup product on the standard complex is functorial. Let
$\varphi:H\to G$ be a group homomorphism, let $M_i$ (respectively, $N_i$)
be $G$-modules (respectively, $H$-modules), and let $f_i:M_i\to N_i$ be a
homomorphism equivariant with respect to $\varphi$, for $i=1,2$; see
Definition~\ref{def:group-equivariance}. Then the following diagram commutes:
% segment-id: o014.li2.chapter6.cup-product.g015
\[\begin{tikzcd}
	C(G, M_1) \otimes C(G, M_2) \arrow[r, "\cup"] \arrow[d] & C(G, M_1 \otimes M_2) \arrow[d] \\
	C(H, N_1) \otimes C(H, N_2) \arrow[r, "\cup"'] & C(H, N_1 \otimes N_2);
\end{tikzcd}\]
The vertical arrows come from the equivariant homomorphisms $f_1$, $f_2$,
and $f_1\otimes f_2$; see Proposition~\ref{prop:grp-equiv}. The normalized
standard complex inherits this property.

% segment-id: o014.li2.chapter6.cup-product.co001
\begin{corollary}\label{prop:cup-ses}
	Suppose there is a short exact sequence of $G$-modules
	$0\to M'_1\xrightarrow{f}M_1\xrightarrow{g}M''_1\to0$ such that the
	corresponding sequence
% segment-id: o014.li2.chapter6.cup-product.q021
	\[ 0 \to M'_1 \otimes M_2 \to M_1 \otimes M_2 \to M''_1 \otimes M_2 \to 0 \]
	is also short exact. Write the connecting homomorphism on cohomology as
	$\delta^p:\Hm^p(G,M''_1)\to\Hm^{p+1}(G,M'_1)$, and so forth. Then
% segment-id: o014.li2.chapter6.cup-product.q022
	\[ (\delta^p \alpha) \cup \beta = \delta^{p+q}(\alpha \cup \beta), \quad \alpha \in \Hm^p(G, M''_1), \; \beta \in \Hm^q(G, M_2). \]

	If instead one takes a short exact sequence
	$0\to M'_2\to M_2\to M''_2\to0$ such that
% segment-id: o014.li2.chapter6.cup-product.q023
	\[ 0 \to M_1 \otimes M'_2 \to M_1 \otimes M_2 \to M_1 \otimes M''_2 \to 0 \]
	remains exact, then
% segment-id: o014.li2.chapter6.cup-product.q024
	\[ \alpha \cup (\delta^q \beta) = (-1)^p \delta^{p+q} (\alpha \cup \beta). \]
\end{corollary}

% segment-id: o014.li2.chapter6.cup-product.pf007
\begin{proof}
	Handle the first case using the cup product on the standard complex. Choose
	a representative $a''\in C^p(G,M''_1)$ of $\alpha$. Recall the concrete
	construction of $\delta^p\alpha$ through the snake lemma: there is an
	$a\in C^p(G,M_1)$ mapping to $a''$, and $da''=0$ implies that $da$ comes
	from some $a'\in C^{p+1}(G,M'_1)$; the cohomology class of the latter
	element is $\delta^p\alpha$. Diagrammatically:
% segment-id: o014.li2.chapter6.cup-product.g016
	\[\begin{tikzcd}
		& a \arrow[mapsto, r] \arrow[mapsto, d, "d"'] & a'' \\
		a' \arrow[mapsto, r] & da &
	\end{tikzcd}\]

	Next, choose a representative $b\in C^q(G,M_2)$ of $\beta$. From $db=0$
	and the Leibniz rule, one obtains $d(a''\cup b)=0$ and
% segment-id: o014.li2.chapter6.cup-product.g017
	\[\begin{tikzcd}
		& a \cup b \arrow[mapsto, r] \arrow[mapsto, d, "d"'] & a'' \cup b \\
		a' \cup b \arrow[mapsto, r] & da \cup b &
	\end{tikzcd}\]
	This implies that the cohomology class of $a'\cup b$ is
	$\delta^{p+q}(\alpha\cup\beta)$.

	The argument for the second case is similar, or it can be reduced to the
	first case by graded commutativity.
\end{proof}

% segment-id: o014.li2.chapter6.cup-product.p025
As an application, we now explain the relationship between the
corestriction $\mathrm{cor}^n$ of
Definition--Proposition~\ref{def:cor} and the cup product.

% segment-id: o014.li2.chapter6.cup-product.pr003
\begin{proposition}[Projection formula]
	Let $M_1$ and $M_2$ be $G$-modules, let $H$ be a subgroup of $G$, and
	suppose $(G:H)$ is finite. For all $\alpha\in\Hm^p(G,M_1)$ and
	$\beta\in\Hm^q(H,M_2)$, the following equality holds in
	$\Hm^{p+q}(G,M_1\otimes M_2)$:
% segment-id: o014.li2.chapter6.cup-product.e004
	\begin{equation*}
		\mathrm{cor}^{p+q}(\mathrm{res}^p(\alpha) \cup \beta) = \alpha \cup \mathrm{cor}^q(\beta);
	\end{equation*}
	and likewise with the left and right positions interchanged.
\end{proposition}

% segment-id: o014.li2.chapter6.cup-product.pf008
\begin{proof}
	First handle the case $p=q=0$. Recall that $\mathrm{cor}^0$ becomes
	$\nu^{G|H}:M_i^H\to M_i^G$, which sends $x\in M_i^H$ to
	$\sum_{gH}gx\in M_i^G$. It therefore suffices to prove that the following
	diagram commutes:
% segment-id: o014.li2.chapter6.cup-product.g018
	\[\begin{tikzcd}
		M_1^G \otimes M_2^H \arrow[rr, "{\identity \otimes \nu^{G|H}}"] \arrow[d] & & M_1^G \otimes M_2^G \arrow[d, "\text{evident}"] \\
		M_1^H \otimes M_2^H \arrow[r, "\text{evident}"'] & (M_1 \otimes M_2)^H \arrow[r, "{\nu^{G|H}}"'] & (M_1 \otimes M_2)^G .
	\end{tikzcd}\]

	This is immediate. If $x\in M_1^G$ and $y\in M_2^H$, then
	$\nu^{G|H}(x\otimes y)=\sum_{gH}gx\otimes gy
	=\sum_{gH}x\otimes gy=x\otimes\nu^{G|H}(y)$.

	For the case $p\geq1$ or $q\geq1$, we use dimension shifting and a
	recursive argument. For example, consider the case $p\geq1$ and the short
	exact sequence of $G$-modules
% segment-id: o014.li2.chapter6.cup-product.q025
	\[ 0 \to M_1 \to \Ind^G_{\{1\}} M_1 \to M'_1 \to 0, \]
	where the first morphism sends $x$ to $[g\mapsto gx]$, for $g\in G$. As a
	homomorphism of $\Bbbk$-modules, this morphism has the left inverse
	$\varphi\mapsto\varphi(1_G)$; hence tensoring the short exact sequence with
	$M_2$ preserves exactness. On the other hand, for $n\geq1$,
	$\Hm^n(G,\Ind^G_{\{1\}}(M_1))\simeq\Hm^n(\{1\},M_1)=0$
	(Theorem~\ref{prop:Shapiro}). Thus the connecting homomorphism gives a
	surjection
% segment-id: o014.li2.chapter6.cup-product.q026
	\[ \delta_G^{p-1}: \Hm^{p-1}(G, M'_1) \twoheadrightarrow \Hm^p(G, M_1). \]

	If this short exact sequence is restricted to $H$, one still obtains a
	connecting homomorphism
	$\delta_H^{p-1}:\Hm^{p-1}(H,M'_1)\to\Hm^p(H,M_1)$, and so forth. Notice
	that $\Res^G_H$ commutes with the operation $\otimes M_2$.

	We may therefore suppose that
	$\alpha=\delta_G^{p-1}(\alpha')$. Compatibility of $\mathrm{res}$ and
	$\mathrm{cor}$ with connecting homomorphisms gives
% segment-id: o014.li2.chapter6.cup-product.g019
	\begin{multline*}
		\mathrm{cor}^{p+q}\left( \mathrm{res}^p(\delta_G^{p-1} \alpha') \cup \beta \right) = \mathrm{cor}^{p+q}\left( \delta_H^{p-1} (\mathrm{res}^{p-1} \alpha') \cup \beta \right) \\
		\xlongequal{\text{Corollary \ref{prop:cup-ses}}} \mathrm{cor}^{p+q} \delta_H^{p+q-1} \left( (\mathrm{res}^{p-1} \alpha') \cup \beta \right) \\
		= \delta^{p+q-1}_G \mathrm{cor}^{p+q-1} \left( (\mathrm{res}^{p-1} \alpha') \cup \beta \right)
		\xlongequal{\text{induction}} \delta^{p+q-1}_G \left(\alpha' \cup \mathrm{cor}^q \beta\right) \\
		\xlongequal{\text{Corollary \ref{prop:cup-ses}}} (\delta^{p-1}_G \alpha') \cup \mathrm{cor}^q \beta .
	\end{multline*}

	The induction on $q$ is entirely similar. The version with the left and
	right positions interchanged follows from graded commutativity.
\end{proof}

% segment-id: o014.li2.chapter6.cup-product.ex001
\begin{example}\label{eg:cup-periodicity}
	Let $C_m$ be a cyclic group of order $m$, where
	$m\in\Z_{\geq1}$. Take the coefficient ring $\Bbbk=\Z$. We explain how to
	understand, through the cup product, the periodic isomorphism
	$\Hm^n(C_m,A)\simeq\Hm^{n+2}(C_m,A)$ implied by
	Proposition~\ref{prop:group-cohomology-cyclic}, where $A$ is an arbitrary
	$C_m$-module and $n\geq1$.

	Choose any generator $t$ of $\Hm^2(C_m,\Z)\simeq\Z/m\Z$. We claim that
	for $n\geq1$, the cup product induces an isomorphism
% segment-id: o014.li2.chapter6.cup-product.q027
	\[ t \cup (\cdot): \Hm^n(C_m, A) \rightiso \Hm^{n+2}(C_m, \Bbbk \otimes A) = \Hm^{n+2}(C_m, A). \]

	Choose a generator $\sigma$ of $C_m$. There are short exact sequences
% segment-id: o014.li2.chapter6.cup-product.tb001
	\[\begin{array}{cc}
		0 \to \mathfrak{I} \to \Z[C_m] \to \Z \to 0, & 0 \to \mathfrak{I} \otimes A \to \Z[C_m] \otimes A \to A \to 0, \\
		0 \to \Z \xrightarrow{\nu} \Z[C_m] \xrightarrow{\sigma - 1} \mathfrak{I} \to 0, & 0 \to A \xrightarrow{\nu \otimes \identity} \Z[C_m] \otimes A \xrightarrow{(\sigma-1) \otimes \identity} \mathfrak{I} \otimes A \to 0.
	\end{array}\]
	The sequences on the left come from
	Lemma~\ref{prop:cyclic-resolution-prep}. Each of their terms is a free
	$\Z$-module, so the sequences on the right obtained by applying
	$\otimes A$ are also exact. First, this gives
% segment-id: o014.li2.chapter6.cup-product.q028
	\[ \Z = \Hm^0(C_m, \Z) \twoheadrightarrow \Hm^1(C_m, \mathfrak{I}) \rightiso \Hm^2(C_m, \Z) = \Z/m\Z, \]
	where both maps are connecting homomorphisms from the long exact
	sequences. Choose a preimage $\tilde{t}$ of $t$ in $\Z$; this integer must
	be relatively prime to $m$. Corollary~\ref{prop:cup-ses} gives the
	commutative diagram
% segment-id: o014.li2.chapter6.cup-product.g020
	\[\begin{tikzcd}[column sep=large]
		\Hm^n(C_m, A) \arrow[r, "\text{connecting homomorphism}"] & \Hm^{n+1}(C_m, \mathfrak{I} \otimes A) \arrow[r, "\text{connecting homomorphism}"] & \Hm^{n+2}(C_m, A). \\
		& \Hm^n(C_m, A) \arrow[lu, "{\tilde{t} \cup (\cdot)}"] \arrow[u] \arrow[ru, "{t \cup (\cdot)}"'] &
	\end{tikzcd}\]

	The exercises in this chapter will show that the composite in the first
	row is the periodic isomorphism. It therefore suffices to prove that
	$\tilde{t}\cup(\cdot)$ is an isomorphism. This is clear: the operation is
	precisely multiplication by $\tilde{t}\in\Z$, while
	$m\Hm^n(C_m,A)=\{0\}$
	(Corollary~\ref{prop:group-coh-torsion}). In fact, if $t$ corresponds to
	$1\bmod m$, take $\tilde{t}=1$; it is then immediately clear that the
	periodic isomorphism is exactly $t\cup(\cdot)$.
\end{example}

% segment-id: o014.li2.chapter6.cup-product.r002
\begin{remark}[Cap product]\label{rem:cap-product}
	\index{cap product}
	As in topology, there is also an operation between group cohomology and
	group homology called the cap product,
% segment-id: o014.li2.chapter6.cup-product.q029
	\[ \cap: \Hm^p(G, M_1) \otimes \Hm_q(G, M_2) \to \Hm_{q-p}(G, M_1 \otimes M_2), \quad M_1, M_2: \text{$G$-modules}. \]
	Since this book does not use the cap product, we mention it only briefly.
	Regard an element $\alpha\in\Hm^p(G,M_1)$ as a morphism
	$\Bbbk\to M_1[p]$ in $\cate{D}(G)$. This element induces the following
	morphism in $\cate{D}^-(G)$:
% segment-id: o014.li2.chapter6.cup-product.g021
	\begin{multline*}
		\Bbbk \otimesL[{\Bbbk[G]}] M_2 \rightiso \Bbbk \otimesL[{\Bbbk[G]}] (\Bbbk \otimesL M_2) \xrightarrow{\identity \otimes \alpha \otimes \identity} \Bbbk \otimesL[{\Bbbk[G]}] (M_1[p] \otimesL M_2) \\
		\xrightarrow{\identity \otimesL \mathrm{can}} \Bbbk \otimesL[{\Bbbk[G]}] (M_1[p] \otimes M_2) \xrightarrow{\identity \otimesL \theta} \Bbbk \otimesL[{\Bbbk[G]}] \left( (M_1 \otimes M_2)[p] \right) \xrightarrow{\theta'} \left(\Bbbk \otimesL[{\Bbbk[G]}] (M_1 \otimes M_2)\right) [p].
	\end{multline*}
	Finally, take $\Hm^{-q}$ of this composite to obtain
% segment-id: o014.li2.chapter6.cup-product.q030
	\[ \alpha \cap (\cdot): \Hm_q(G, M_2) \to \Hm_{q-p}(G, M_1 \otimes M_2). \]
	In the special case $p=q$, one obtains a pairing
% segment-id: o014.li2.chapter6.cup-product.q031
	\[ \lrangle{\cdot, \cdot}: \Hm^p(G, M_1) \otimes \Hm_p(G, M_2) \to \left( M_1 \otimes M_2 \right)_G . \]
\end{remark}
